\chapter{Conclusion}

Much can be done to help preserve the history of modeling and simulation in Austin and Houston from 1958 to 1982, the era of the mainframes, focusing on the code and the stories of the people who created it. They had vision enough, using very primitive hardware, to be on the frontier opening up a new field that essentially didn’t exist before the electronic digital computer. It’s probably enlightening to start with the human computers at Los Alamos and their rapid adoption of the forerunner of the IBM CPC, follow that thread to Santa Monica and Houston, then onward to Austin, looking at a handful of real breakthroughs along the way. Simplex, SOR, ADI, subsurface, and orbit determination. Possibly it’s worthwhile to trace back to earlier mechanical, analog modeling of reservoirs and wind tunnels, electro-mechanical computers, etc. It’s worth noting that these developments were not foreseen and not popularly recognized at the time or since. Science fiction and popular science promoted space travel and artificial brains, but not the somewhat invisible and mathematical world of modeling and simulation.

Above all these models were linear algebra systems, and that’s crucial from a historical and preservation viewpoint. We might not have original source code for subsurface modeling from 1953, but we know the machines, the IBM CPC for example, and we know the exact mathematical methods involved. Is there an important distinction between having the original source code and reconstructing it using the same machines, tools, and, most importantly, the same linear algebra techniques? The assembly and Fortran will be effectively the same in any case, and we can absolutely reconstruct and run the models of the early days. The code was concise in modern terms, but powerful, and we can preserve the original purity of key models that changed the course of industries and scientific fields. 

There are working, usable examples of robust time capsule docker containers, each a self-contained environment dedicated to a model or simulation with its hardware, operating system, tools, code, and documentation all version controlled and fully automated locally and in the cloud. Docker emphasises textuality, with all of the artifacts and materials in readable form, and we are just beginning to explore how AI can play a critical role in this environment, allowing us to do more work, with more effective power, than would be possible otherwise.

Beyond a historical archive for research purposes, these docker time capsules can also serve as a library for modelers from industry, engineering, and the physical sciences, and there’s a historical parallel to draw on to motivate that. Broadly speaking, the types of people concerned will naturally have an inherent interest in physics and will be aware of The Feynman Lectures on Physics and its special place in the field. These lectures are a role model and path finder for the type of time capsules contemplated here. One of the things that Feynman does is consistently bring matters back to first-principles, clarity, and a certain minimalist purity. If these time capsules could someday be praised as “a kind of Feynman Lectures of modeling” then they will have succeeded in meeting their potential. 

A famous piece from the lectures, table 22-3 from the lecture on algebra, is adopted as the simplest model to demonstrate the health of a time capsule. What makes it remarkable is that it links two entirely different worlds, algebra and geometry, using nothing more than basic addition and multiplication. Without touching trigonometry or geometry, pure arithmetic organically produces smooth, continuous oscillation. The continuous motion of waves and orbits, all the rhythmic, circular phenomena in the physical universe, are naturally encoded inside the basic machinery of exponents and imaginary numbers. The few lines of assembly or Fortran to generate this table are implementable in any healthy time capsule and can act as a kind of heartbeat and a fundamental sign of life.
